%=============================================================
% Paper 4: Davis & Haltiwanger (1999), AER
%=============================================================
\chapter{Davis \& Haltiwanger (1999)}

\begin{paperbox}{On the Driving Forces Behind Cyclical Movements in Employment
and Job Reallocation}
\textbf{Authors:} Steven J.\ Davis \& John Haltiwanger\\[3pt]
\textbf{Journal:} \textit{American Economic Review} 89(5): 1234--1258\\[3pt]
\textbf{Year:} 1999 \quad\textbf{Field:} Macroeconomics; Labour economics\\[3pt]
\textbf{Classification:} Empirical (structural VAR)
\end{paperbox}

%------------------------------------------------------------
\section{Research Question}

What types of structural disturbances---allocative shocks (which shift the
distribution of employment opportunities across job sites) versus aggregate
shocks (which move all plants in the same direction)---are the primary drivers
of cyclical movements in aggregate employment and job reallocation in U.S.\
manufacturing?

%------------------------------------------------------------
\section{Motivation}

The DH (1992) paper documented that job reallocation is countercyclical and
job destruction substantially more volatile than job creation. A key open
question is whether these patterns reflect primarily allocative disturbances
(sector-specific technology, oil price, or structural change shocks) or
aggregate disturbances (monetary policy, aggregate demand, aggregate
productivity). This distinction matters enormously for policy: if reallocation
is primarily driven by allocative shocks, standard aggregate demand
stabilisation cannot reduce it and might not even be desirable; if aggregate
shocks drive most employment fluctuations, conventional macroeconomic
interventions are more directly relevant. Previous VAR identification
strategies (Blanchard--Quah 1989; Bernanke 1986) could not exploit the
information in gross job flows. Davis--Haltiwanger (1999) develops a novel
identification strategy using the creation--destruction decomposition that
allows both qualitative and quantitative restrictions motivated by economic
theory.

%------------------------------------------------------------
\section{Main Contribution}

\begin{enumerate}
  \item \textbf{Identification via creation--destruction:} Develops a
    structural VAR where the decomposition of employment changes into creation
    and destruction rates provides novel identifying information unavailable
    in aggregate employment data alone. Aggregate and allocative disturbances
    have \emph{qualitatively different} implications for creation and
    destruction: aggregate shocks move them in \emph{opposite} directions,
    while allocative shocks move them in the \emph{same} direction.

  \item \textbf{Tighter inequality restrictions:} Theory implies that job
    destruction responds \emph{at least as much} as job creation to aggregate
    shocks ($b_{na}\leq -1$), and that the contemporaneous impact of allocative
    shocks on job creation does not exceed that on job destruction ($|b_{ps}|\leq 1$).
    These restrictions, drawn from search-theoretic models and option-value
    arguments, substantially narrow the range of permissible inferences.

  \item \textbf{Long historical series:} Estimates the system over
    1947:Q1--1993:Q4, exploiting a long quarterly series constructed by
    splicing LRD data (1972--1993) with BLS turnover data (1947--1981) using
    a cyclically varying quit replacement rate.
\end{enumerate}

%------------------------------------------------------------
\section{Relationship to the Literature}

The paper is a direct sequel to DH1992 and DH1990. It is methodologically
related to the structural VAR literature (Blanchard--Quah 1989; King--Watson
1992; Shapiro--Watson 1988) but extends it by using the creation--destruction
decomposition to generate additional identifying restrictions not available in
aggregate time series. It contributes to the debate between ``cleansing
recession'' theories (Caballero--Hammour 1994; Lilien 1982) and ``sullying''
theories, and shows that the identification of which shock dominates depends
critically on which restrictions one is willing to impose.

%------------------------------------------------------------
\section{Theoretical Framework and Economic Mechanism}

\subsection*{Structural Moving Average Representation}

Let $Y_t=[\text{POS}_t,\,\text{NEG}_t]'$ and
$\varepsilon_t=[\varepsilon^a_t,\,\varepsilon^s_t]'$ (aggregate and
allocative structural innovations). The structural representation is:
\begin{equation}
  Y_t = B(L)\,\varepsilon_t, \qquad B(0) = B_0.
  \label{eq:structural_ma}
\end{equation}
The estimated reduced-form VAR yields:
\begin{equation}
  Y_t = D(L)\,\eta_t, \qquad D(0) = I,
  \label{eq:rf_var}
\end{equation}
with reduced-form innovations $\eta_t = B_0\,\varepsilon_t$.

\subsection*{Identification Problem}

With diagonal elements of $B_0$ normalised to unity, there are two free
parameters: $b_{na}$ (contemporaneous effect of an aggregate shock on POS
relative to NEG) and $b_{ps}$ (contemporaneous effect of an allocative shock
on NEG relative to POS). The zero-covariance restriction
$\text{Cov}(\varepsilon^a,\varepsilon^s)=0$ yields the moment equation:
\begin{equation}
  b_{ps} = \frac{\sigma_{pn} - b_{na}\,\sigma_p^2}{\sigma_n^2 - b_{na}\,\sigma_{pn}},
  \label{eq:dh99_bps}
\end{equation}
leaving the system one-dimensionally underidentified.

\subsection*{Weak Inequality Restrictions}

\begin{align}
  &\text{(i)}\quad b_{na} < 0:
    &&\text{aggregate shocks move POS and NEG in opposite directions}\\
  &\text{(ii)}\quad b_{ps} > 0:
    &&\text{allocative shocks move both in the same direction}
\end{align}

\subsection*{Tighter Theory-Motivated Restrictions}

\begin{align}
  &\text{(i)'}\quad b_{na} \leq -1:
    &&\text{NEG responds at least as much as POS to aggregate shocks}\\
  &\text{(ii.a)'}\quad |b_{ps}| \leq 1:
    &&\text{allocative shock impact on POS}\leq\text{impact on NEG}\\
  &\text{(ii.b)'}\quad \textstyle\sum_{l=1}^{m} B_{11}(l) > 0
    &&\text{for all } m \leq M\!=\!16
\end{align}

Restriction (i)' follows from search models: separations can occur
instantaneously while new matches require search time, so an aggregate shock
causes a sharper contemporaneous rise in destruction than fall in creation.
Restriction (ii.a)' follows from option value arguments: sunk costs of
investment create an incentive to delay job creation in the face of uncertain
future allocative disturbances, depressing creation even as destruction rises.

\subsection*{Long-Run Neutrality Restrictions (Alternative)}

\begin{align}
  &\text{(iv)}\quad \sum_{l=0}^{\infty}\bigl[B_{11}(l) + B_{21}(l)\bigr] = 0:
    &&\text{aggregate shocks have no long-run effect on job reallocation}\\
  &\text{(v)}\quad \sum_{l=0}^{\infty}\bigl[B_{12}(l) - B_{22}(l)\bigr] = 0:
    &&\text{allocative shocks have no long-run effect on employment}
\end{align}

%------------------------------------------------------------
\section{Data}

\begin{itemize}
  \item \textbf{Primary (1972--1993):} LRD quarterly job creation and
    destruction data, updated through 1993 (Davis et al.\ 1996).
  \item \textbf{Secondary (1947--1981):} BLS monthly turnover data (accessions,
    layoffs, quits), converted using the Blanchard--Diamond (1990) method with
    a \emph{cyclically varying} quit replacement rate.
  \item \textbf{Splicing:} Join the two series over their 1972--1981 overlap.
    $\rho(\text{implied NET growth},\text{BLS 790 data})=0.91$.
  \item \textbf{Sample:} 1947:Q1--1993:Q4 (187 observations).
\end{itemize}

\begin{center}
\begin{tabular}{lcccc}
\toprule
\textbf{Statistic} & \textbf{POS} & \textbf{NEG} & \textbf{NET} & \textbf{SUM} \\
\midrule
Mean & 5.8\% & 6.0\% & $-0.1\%$ & 11.8\% \\
Std dev & 1.2\% & 1.5\% & 2.0\% & 1.7\% \\
Min & 3.8\% & 2.9\% & $-5.9\%$ & 7.4\% \\
Max & 10.2\% & 10.8\% & 6.8\% & 16.3\% \\
\bottomrule
\multicolumn{5}{l}{\small Source: BLS--LRD spliced series, 1947:Q1--1993:Q4}
\end{tabular}
\end{center}

%------------------------------------------------------------
\section{Empirical Strategy}

Structural VAR estimation with multiple identification schemes:
\begin{enumerate}
  \item Estimate a 2-variable VAR in POS and NEG with 4 lags (Dickey--Fuller
    tests reject unit roots in both series).
  \item Derive the reduced-form covariance matrix.
  \item Apply each set of identifying restrictions to recover admissible ranges
    for $(b_{na}, b_{ps})$ and corresponding variance decompositions.
  \item Bootstrap standard errors using 1,000 Monte Carlo draws.
  \item Construct historical decompositions at boundary values of the
    admissible parameter space (Figures 4--5).
\end{enumerate}

%------------------------------------------------------------
\section{Identification Strategy}

This is the paper's central methodological innovation. Rather than point-
identifying the system, the paper derives \emph{ranges} of parameter values
consistent with a set of qualitative restrictions, then reports the implied
range of variance decompositions.

Key identification elements:
\begin{itemize}
  \item \textbf{Sign restrictions on contemporaneous responses:} (i) and (ii)
    are definitional and accepted by nearly all theories.
  \item \textbf{Magnitude restrictions:} (i)' and (ii.a)' are theory-motivated
    from search models and option value arguments; these are the \emph{binding}
    restrictions that narrow the admissible range.
  \item \textbf{Zero covariance:} Standard in the VAR literature; the paper
    explores sensitivity to relaxing this (Figures 7--8).
\end{itemize}

\textbf{Critical limitation:} The mapping from $(b_{na}, b_{ps})$ to variance
decompositions is highly nonlinear (Figure~3). The allocation of variance
between aggregate and allocative shocks is extremely sensitive to where in the
admissible range one sits. The paper honestly reports ranges rather than point
estimates---an intellectually rigorous but informationally limited strategy.

%------------------------------------------------------------
\section{Main Results}

\subsection*{Under Weak Restrictions (Figure~3)}

Admissible range: $b_{na}\in(-2.5,-0.87)$.
\begin{itemize}
  \item Allocative shocks account for 5--20\% of 4-step forecast-error
    variance of employment growth.
  \item Aggregate shocks are always a major driving force behind employment
    fluctuations.
  \item Allocative shocks account for 20--90\% of job reallocation variance
    (wide range).
\end{itemize}

\subsection*{Under Tighter Restrictions (Table~3, Columns 1--2)}

Admissible range narrows to $b_{na}\in(-1.63,-1.0)$, $b_{ps}\in(0.1,1.0)$.
\begin{center}
\begin{tabular}{lcc}
\toprule
\textbf{Variable} & \textbf{Range (4-step horizon)} &
  \textbf{Range (16-step horizon)} \\
\midrule
Employment growth & 3\%--13\% & 5\%--17\% \\
Job reallocation & 44\%--80\% & 49\%--84\% \\
\bottomrule
\multicolumn{3}{l}{\small Fraction of forecast-error variance explained by
\emph{allocative} shocks.}
\end{tabular}
\end{center}

\textbf{Key inference:} Allocative disturbances are a \emph{minor} driver of
employment fluctuations but the \emph{dominant} driver of job reallocation
intensity.

\subsection*{Under Neutrality Restrictions (Table~3, Columns 3--5)}

\begin{itemize}
  \item Restricting allocative shocks to have no long-run employment effect
    (v) yields $b_{na}=0.03$, $b_{ps}=-0.41$---\emph{violates} weak
    qualitative restrictions (i), (ii), (i)'.
  \item Restricting aggregate shocks to have no long-run reallocation effect
    (iv) yields $b_{na}=-0.68$, $b_{ps}=-0.12$---violates (i)' and (i).
  \item Under these neutrality restrictions, allocative shocks account for
    34--80\% of employment variance and $>$80--94\% of reallocation variance.
\end{itemize}

%------------------------------------------------------------
\section{Economic Magnitude and Interpretation}

The paper's core quantitative finding under the tighter qualitative restrictions
is:
\begin{itemize}
  \item \textbf{Aggregate shocks account for 83--97\%} of cyclical variation
    in employment growth (4-step forecast errors).
  \item \textbf{Allocative shocks account for 44--80\%} of cyclical variation
    in job reallocation intensity.
\end{itemize}
These two findings are not in tension: aggregate shocks cause recessions;
allocative shocks cause the gross flows that accompany them independently.
The economy's output cycles and its reallocation cycles are driven by different
forces.

%------------------------------------------------------------
\section{Mechanisms}

Two mechanisms determine the relative importance of the two shock types:

\begin{enumerate}
  \item \textbf{Opportunity cost mechanism} (Davis--Haltiwanger 1990): during
    recessions, the opportunity cost of relocating workers falls, so the pace
    of allocative reallocation endogenously increases. A constant stream of
    allocative disturbances generates countercyclical reallocation intensity.

  \item \textbf{Option value / sunk cost mechanism:} The sunk nature of
    investments required to create new vacancies generates an option value
    for waiting. When allocative shocks are serially correlated (Table~2
    shows significant autocorrelations in cross-industry return dispersion
    at lags 1--6), this option value simultaneously depresses job creation
    and boosts job destruction---enabling $b_{ps}<0$, i.e., allocative
    innovations that cause short-run employment declines.
\end{enumerate}

%------------------------------------------------------------
\section{Robustness Checks}

\begin{itemize}
  \item Results reported for five different identification schemes (qualitative
    and neutrality restrictions).
  \item Sensitivity analysis for $\rho(\varepsilon^a,\varepsilon^s)
    \in\{-0.3,\,0,\,+0.3\}$ (Figures 7--8).
  \item BLS~790 data cross-validation ($\rho = 0.91$ between implied and
    published net employment growth).
  \item Dickey--Fuller tests confirm stationarity of POS and NEG series.
  \item VARs with deterministic trends as sensitivity checks.
\end{itemize}

%------------------------------------------------------------
\section{Limitations}

\begin{enumerate}
  \item \textbf{One-dimensional underidentification:} No single restriction
    produces a unique parameter point; results are necessarily range-valued.
  \item \textbf{Two-variable system:} Cannot distinguish among multiple types
    of allocative or aggregate shocks (technology, monetary, fiscal, oil).
  \item \textbf{Manufacturing only:} The long historical series is available
    only for manufacturing, which exhibits secular decline over the sample.
  \item \textbf{BLS splicing:} The 1947--1971 portion relies on a maintained
    assumption about quit replacement rates, introducing potential measurement
    error.
  \item \textbf{Linearity:} Nonlinear interactions between aggregate and
    allocative shocks---which the theories all suggest---are not accommodated.
  \item \textbf{Inconsistency of neutrality and qualitative restrictions:}
    Long-run neutrality restrictions produce $(b_{na}, b_{ps})$ values that
    violate the sign restrictions, revealing internal inconsistency and
    limiting the credibility of the long-run results.
\end{enumerate}

%------------------------------------------------------------
\section{Policy Implications}

\begin{itemize}
  \item Because aggregate shocks drive most cyclical employment fluctuations
    (under preferred tighter restrictions), standard aggregate demand policies
    can influence employment---but do not address the underlying allocative
    process that determines job reallocation intensity.
  \item The finding that allocative shocks are a minor driver of employment
    cycles but a dominant driver of reallocation suggests that policies
    targeting ``structural unemployment'' address a phenomenon largely
    independent of macroeconomic stabilisation.
  \item The long-run neutrality results---where allocative shocks account for
    34--80\% of employment variance---would, if taken seriously, suggest a
    much larger role for structural policies. However, these results are
    theoretically inconsistent with the qualitative restrictions.
\end{itemize}

%------------------------------------------------------------
\section{Overall Assessment}

Davis--Haltiwanger (1999) is a methodologically important paper that introduces
a genuinely novel approach to structural VAR identification using the
creation--destruction decomposition. It is honest about the limits of what the
data can identify and clearly maps the trade-offs between different sets of
assumptions. The main empirical finding---under theory-motivated tighter
restrictions, aggregate shocks drive most employment cyclicality while
allocative shocks drive most reallocation cyclicality---is sensible and has
influenced the subsequent literature. The paper's weakness is that its
conclusions are range-valued rather than point estimates, and the ranges are
wide enough to be somewhat uninformative on policy-relevant questions.

\textbf{Quality: Very good. A methodologically careful paper with important
empirical findings, though limited by data and identification constraints.}

%------------------------------------------------------------
\section{Relevance to My Research}

The identification strategy---using the creation--destruction decomposition to
distinguish allocative from aggregate shocks---is directly applicable to
Canadian data. With LEAP/LEED, one could construct quarterly job creation and
destruction series for the entire Canadian economy and estimate the DH1999
structural VAR. The Canadian context is particularly interesting because the
economy faces both aggregate shocks (U.S.\ monetary policy spillovers, global
demand) and strong allocative shocks (commodity price cycles with large
sector-specific effects). Testing whether the relative importance of allocative
vs.\ aggregate shocks differs in resource-intensive provinces versus
manufacturing provinces would be a natural extension.

%------------------------------------------------------------
\section{Potential Research Extensions}

\begin{itemize}
  \item Extend the DH1999 VAR to the service sector in Canada using the full
    LEAP/LEED universe, which would make the sample more representative of
    the aggregate economy.
  \item Add commodity price innovations and exchange rate movements as
    additional variables in an expanded VAR, exploiting Canadian resource-sector
    dynamics to better identify allocative shock contributions.
  \item Use regional variation in commodity price exposure to construct
    province-specific identifying restrictions on allocative shock magnitudes.
\end{itemize}

%------------------------------------------------------------
\section{Research Gaps}

\begin{itemize}
  \item The paper uses only two variables (POS, NEG). Expanding to a
    four-variable system (POS, NEG, unemployment rate, vacancy rate) would
    permit Beveridge curve restrictions and substantially tighten identification.
  \item The allocation between aggregate and allocative shocks is sensitive to
    $\rho(\varepsilon^a,\varepsilon^s)$ but the paper has no empirical estimate
    of this parameter.
\end{itemize}

%------------------------------------------------------------
\section{Methodological Lessons}

\begin{enumerate}
  \item The creation--destruction decomposition provides identifying information
    unavailable in aggregate net employment data: allocative and aggregate
    shocks have qualitatively different creation--destruction responses, and
    theory imposes magnitude restrictions.
  \item When a system is underidentified, reporting the range of inferences
    over the admissible parameter space (rather than arbitrary point
    identification) is the intellectually honest approach.
  \item Long-run neutrality restrictions and short-run qualitative restrictions
    can be internally inconsistent (Table~3), providing a diagnostic for
    whether theoretical identifying assumptions are self-consistent.
  \item The serial correlation of cross-industry stock return dispersion
    (Table~2) as indirect evidence for the option value mechanism is a simple
    and elegant test of a theoretical underpinning of an identifying assumption.
\end{enumerate}

%------------------------------------------------------------
\section*{Five-Sentence Summary}

Davis and Haltiwanger (1999) develop a structural VAR framework for a
bivariate system of quarterly job creation and destruction rates in U.S.\
manufacturing (1947--1993), exploiting the qualitatively different responses
of creation and destruction to aggregate versus allocative shocks to generate
novel identifying restrictions unavailable in aggregate employment data alone.
Weak inequality restrictions establish that aggregate shocks are always major
driving forces behind employment fluctuations, while the contribution of
allocative shocks to employment variance ranges widely (5--20\%) depending on
the structural parameter values. Imposing tighter theory-motivated
restrictions---that job destruction responds at least as much as creation to
aggregate shocks, and that allocative shocks cannot simultaneously boost
employment---narrows the range to 3--17\% of employment variance explained by
allocative shocks but confirms that allocative shocks account for 44--80\% of
job reallocation intensity variance. Long-run neutrality restrictions imply an
even larger role for allocative shocks in both employment and reallocation
variance, but these restrictions are internally inconsistent with the
qualitative sign restrictions. The paper's core message is that aggregate
shocks drive most cyclical employment fluctuations, while allocative
disturbances are the dominant driver of the countercyclical movements in job
reallocation that accompany those cycles.

\vspace{6pt}
\begin{quickref}
\begin{tabularx}{\linewidth}{@{}lX@{}}
\textbf{Key contribution} &
  Novel structural VAR identification using job creation--destruction
  decomposition; maps the full range of inferences over the admissible
  structural parameter space under qualitative restrictions from theory. \\[4pt]
\textbf{Key weakness} &
  System remains one-dimensionally underidentified; wide admissible ranges
  limit sharp empirical conclusions; manufacturing-only sample; long-run
  and short-run restrictions are internally inconsistent. \\[4pt]
\textbf{Most interesting gap} &
  Expanding to a four-variable system (POS, NEG, unemployment, vacancies)
  with Beveridge curve restrictions would tighten identification further and
  address whether the allocative--aggregate decomposition is stable across
  business cycle phases. \\[4pt]
\textbf{Publishable extension} &
  Construct a DH1999-type structural VAR for the full Canadian economy using
  quarterly LEAP/LEED-based creation and destruction series, with an expanded
  variable set including commodity price indices to separately identify
  resource-sector allocative shocks.
\end{tabularx}
\end{quickref}
